% Appendix B -- Glossary (PUB-REF-0001 AN-23, F-IBM: term in the side column,
% definition in the text column).  Every \term{} of this book is here
% (CP-03); definitions follow the behaviour documented in the chapters.
% GATES: test-pubs-build
\Gates{test-pubs-build}
\chapter{Glossary}

\SideHead{archive attribute}A mark on a file showing that it has been
created or changed since it was last backed up.  \cmd{ATTRIB} displays it
and \cmd{XCOPY} uses it.

\SideHead{batch file}A text file of commands, with the extension
\file{.BAT}, that InitechDOS carries out when you type its name.

\SideHead{byte}The unit in which file sizes and disk space are measured.  A
byte holds one character of text.

\SideHead{command}An instruction you type at the command prompt.

\SideHead{command processor}The part of InitechDOS, \file{COMMAND.COM},
that displays the prompt and carries out commands.

\SideHead{configuration command}A line in \file{CONFIG.SYS} that sets up
the system when it starts.

\SideHead{current directory}The directory in which commands look for files
unless you name another.  The command prompt shows it.

\SideHead{directory}A list of files on a disk, kept under a name.

\SideHead{diskette}A flexible disk in a protective jacket, used in a
diskette drive.

\SideHead{environment}An area of memory holding named settings, the
environment variables, for InitechDOS and for programs.

\SideHead{environment variable}A named setting in the environment, such as
\texttt{PATH}.

\SideHead{extension}The part of a file name after the period: up to three
characters.

\SideHead{external command}A command that is a program file or batch file
on disk.

\SideHead{file}A named collection of information stored on a disk.

\SideHead{filter}A program that reads its standard input, changes it, and
writes the result to its standard output.

\SideHead{fixed disk}A disk built into the computer.

\SideHead{formatted}Prepared for use by InitechDOS.  A disk must be
formatted before files can be stored on it.

\SideHead{internal command}A command built into the command processor.

\SideHead{lost allocation units}Space on a disk marked as in use that
belongs to no file.  \cmd{CHKDSK} finds it.

\SideHead{operating system}The program that controls the computer, its
disks and devices, and runs other programs.

\SideHead{parameter}Information typed after a command name that tells the
command what to work on.

\SideHead{parent}The directory that contains a subdirectory, named
\file{..}.

\SideHead{path}The list of directory names, separated by backslashes, that
leads to a file or directory.

\SideHead{pipe}The character \texttt{|} between two commands, which sends
the output of the first to the second as its input.

\SideHead{read-only}A file that can be read and copied but not changed or
deleted.

\SideHead{redirection}Sending a command's output to a file or device, or
taking its input from one, with \texttt{>}, \texttt{>>} or \texttt{<}.

\SideHead{root directory}The main directory of a disk, created when the
disk is formatted.

\SideHead{standard input}Where a command reads its input: the keyboard,
unless redirected.

\SideHead{standard output}Where a command writes its output: the display,
unless redirected.

\SideHead{subdirectory}A directory inside another directory.

\SideHead{template}The last line typed, kept so that the editing keys can
use it to build the next line.

\SideHead{tree}The arrangement of the directories of a disk, with the root
directory at its base.

\SideHead{volume label}A name of up to eleven characters given to a disk.

\SideHead{wildcard}The characters \texttt{*} and \texttt{?}, used in a file
name to name a group of files.
